% ARCHIVED DRAFT: independently audited by root and incorporated on 2026-11-19 16:42 UTC.
% The assembled paper is authoritative; this file retains drafting provenance.
\begin{saresult}{B.072}{Proved}{Local bases whose maximal ideal is killed by Frobenius}
Let $(R,\mathfrak m)$ be a nonzero local $\mathbb F_p$-algebra with
finite residue field $\mathbb F_q$. Suppose that $x^p=0$ for every
$x\in\mathfrak m$; the ideal $\mathfrak m$ need not be nilpotent or
finitely generated. For every unital inclusion $R\subset S$, the following
are equivalent:
\[
\begin{gathered}
 \text{the polynomial co-operations exist uniquely};\\
 \theta_2\text{ is injective};\\
 R\longrightarrow S\text{ is a ring epimorphism and }
                     (R:_S\mathfrak m)=R.                \tag{B72.1}
\end{gathered}
\]
Every positive case has all multivariable polynomial tensor isomorphisms.
Thus the nonflat example B.067 lies in a fully classified class of
possibly infinite, non-Noetherian bases. If $\mathfrak m$ is nilpotent as
an ideal, these conditions force $S=R$.
\end{saresult}
\begin{proof}
The image $R^p$ of Frobenius is a coefficient field $k\simeq\mathbb F_q$.
Indeed the kernel of Frobenius is exactly $\mathfrak m$: it contains
$\mathfrak m$ by hypothesis and contains no unit. The image is a field
and maps isomorphically onto the residue field, since Frobenius on that
finite field is an automorphism. Every element of $R$ has a unique
expression $a+n$ with $a\in k$, $n\in\mathfrak m$.

Put
\[
                    Q(X)=(X^q-X)^p=X^{qp}-X^p.
\]
It is monic and vanishes on $R$. The Chinese remainder isomorphism
\[
 k[X]/Q\ \simeq\prod_{a\in k}k[t]/(t^p),
 \qquad X\longmapsto a+t                              \tag{B72.2}
\]
gives a basis $h_{a,j}$, $a\in k$, $0\le j<p$, of polynomials of degree
less than $qp$ whose image in the $a$-component is $t^j$ and whose other
components vanish. Monic division by $Q$ and this basis work over every
$k$-algebra $S$. For $n\in\mathfrak m$ one has
$h_{a,j}(b+n)=0$ if $a\ne b$, and $h_{a,j}(a+n)=n^j$.

Write $C_j=(R:_S\mathfrak m^j)$, so $C_0=R$.
For a low-degree remainder $H=\sum c_{a,j}h_{a,j}$, preservation of $R$
is exactly the condition $c_{a,j}\in C_j$. To verify necessity, evaluate
$H(a+\lambda n)$ for fixed $a,n$ and for $p$ distinct $\lambda\in k$.
The Vandermonde matrix is invertible over $k$, so each $c_{a,j}n^j$ lies
in $R$. For $1\le j<p$, the ideal generated by all $n^j$, $n\in\mathfrak m$,
is $\mathfrak m^j$: the polarization identity
\[
 j!\,n_1\cdots n_j=
 \sum_{J\subset\{1,\ldots,j\}}(-1)^{j-|J|}
                         \left(\sum_{i\in J}n_i\right)^j
\]
and invertibility of $j!$ prove the reverse inclusion. The case $j=0$
is ordinary evaluation. Sufficiency follows by direct substitution.
Hence, as an explicitly split $R$-module,
\[
 \Omega(R,S)=\bigoplus_{a\in k,\ 0\le j<p} C_jh_{a,j}
                                \ \oplus\ Q(X)S[X].       \tag{B72.3}
\]

Assume $\theta_2$ is injective. The last summand of (B72.3) makes the
multiplication map $S\otimes_R S\to S$ injective, and it is already
surjective; hence $R\to S$ is an epimorphism. The summand indexed by
$(a,j)=(0,1)$ similarly makes $C_1\otimes_R C_1\to S$ injective.
Here the leading nonzero coefficient of $h_{0,1}$ belongs to $k^\times$,
so a coefficient tensor has zero polynomial image precisely when its
product in $S$ is zero.

The module $C_1$ is in fact a subring of $S$, with
$\mathfrak mC_1=\mathfrak m$. If $x\in\mathfrak m$ and $c\in C_1$,
then $xc\in R$ and $(xc)^p=0$; thus $xc\in\mathfrak m$.
This proves the ideal equality and shows closure of $C_1$ under products:
$\mathfrak m C_1C_1\subset\mathfrak mC_1\subset R$.
Consequently multiplication $C_1\otimes_R C_1\to C_1$ is both injective
and surjective, making $R\to C_1$ an epimorphism.
Its reduction is an injective epimorphism from the field
$R/\mathfrak m$ to $C_1/\mathfrak m$. A field admits no proper injective
epimorphic algebra: the tensors $c\otimes1$ and $1\otimes c$ are distinct
whenever $1,c$ are linearly independent. Therefore
$C_1/\mathfrak m=R/\mathfrak m$, and $C_1=R$.
This proves necessity in (B72.1).

Suppose conversely that $R\to S$ is an epimorphism and $C_1=R$.
Iterating colons gives
\[
 (R:_S\mathfrak m^{j+1})
       =((R:_S\mathfrak m^j):_S\mathfrak m)=R
                         \qquad(j\ge0).                  \tag{B72.4}
\]
Thus (B72.3) simplifies to
$\Omega(R,S)=R[X]_{<qp}\oplus Q(X)S[X]$.
In several variables, divide by $Q$ in each variable. Every basis term
with at least one $Q$ factor is a null polynomial and has an arbitrary
coefficient in $S$. For the remaining terms use the product Hermite basis.
Independent scalar Vandermonde interpolation in each nilpotent increment
forces a coefficient with jet indices $j_1,\ldots,j_n$ to multiply
$\mathfrak m^{j_1+\cdots+j_n}$ into $R$; polarization in each coordinate
justifies this whole ideal, not only individual powers. Equation (B72.4)
puts that coefficient in $R$.

The tensor power of the one-variable decomposition has the same
coefficient modules: $R$ on every all-low-degree term, and $S$ on every
term with a high-degree factor, since $S\otimes_R S=S$.
The monic bases $X^i$ for $i<qp$ and $Q(X)X^{i-qp}$ for $i\ge qp$
therefore identify every polynomial tensor map with these coefficient
isomorphisms. All $\theta_n$ are isomorphisms.
For explicit co-operations on the new generators, one can also use
\[
\begin{aligned}
 Q(X+Y)&=Q(X)+Q(Y),\\
 Q(XY)&=Q(X)Q(Y)+Q(X)Y^p+X^pQ(Y).
\end{aligned}
\]
These identities put the coaddition and comultiplication of every $cQ$
in the tensor image. Together with ordinary polynomials they generate
the whole analyser, and uniqueness and compatibility follow from the
tensor isomorphisms.

Finally, if $\mathfrak m^N=0$, equation (B72.4) gives
$S=(R:_S0)=R$. This proves the last assertion and also distinguishes
elementwise Frobenius nilpotence from nilpotence of the ideal.
\end{proof}
